% Generated draft: requires mathematical and source review.
\documentclass[11pt]{article}
\usepackage[margin=25mm]{geometry}
\usepackage{fontspec}
\setmainfont{Latin Modern Roman}
\usepackage{amsmath,amssymb,mathtools,unicode-math}
\setmathfont{Latin Modern Math}
\usepackage{enumitem,needspace,xurl,hyperref}
\setlength{\parindent}{0pt}
\setlength{\parskip}{5pt}
\setlength{\emergencystretch}{3em}
\newcommand{\N}{\mathbb{N}}
\newcommand{\Z}{\mathbb{Z}}
\newcommand{\Q}{\mathbb{Q}}
\newcommand{\R}{\mathbb{R}}
\newcommand{\C}{\mathbb{C}}
\newcommand{\F}{\mathbb{F}}
\newcommand{\A}{\mathbb{A}}
\newcommand{\PP}{\mathbb{P}}
\newcommand{\E}{\mathbb{E}}
\newcommand{\eps}{\varepsilon}
\newcommand{\dd}{\,\mathrm{d}}
\begin{document}
\section*{UnsolvedMath ID: 30001783}
\textbf{Source problem code:} catalogue entry
\subsection{Definitions and mathematical statement}
The supplied question concerns translation surfaces but does not specify compactness, genus, stratum, or the ambient space in which divergence is measured. The following gives a precise conditional formulation in the customary compact setting; these conventions require confirmation against the source.

Let \(X=(S,\omega)\) be a compact connected Riemann surface with a nonzero holomorphic one-form, normalized so that its flat area measure \(\mu_X\) satisfies \(\mu_X(S)=1\). Away from the zeros of \(\omega\), integration of \(\omega\) supplies Euclidean coordinates whose transition maps are translations. Let \(\Theta=\R/\pi\Z\) be the space of unoriented directions, equipped with its angular metric. For \(\theta\in\Theta\), let \(\varphi_t^\theta\) denote straight-line flow in direction \(\theta\), defined for almost every point for all times. Reversing the orientation does not change ergodicity.

Define the set of non-ergodic directions, with respect to flat area, by
\[
N(X)=\left\{\theta\in\Theta:\ \exists\text{ a measurable }A\subseteq S,\quad 0<\mu_X(A)<1,\quad \mu_X\bigl((\varphi_t^\theta)^{-1}(A)\mathbin{\triangle}A\bigr)=0\ \text{for every }t\in\R\right\}.
\]
Here \(\triangle\) denotes symmetric difference. This definition concerns failure of ergodicity of flat area; failure of unique ergodicity is a different possible intended meaning and is not substituted for it.

To formulate divergence, choose an ambient moduli space \(\mathcal M\) containing the relevant area-one translation surfaces, with its usual topology. Write
\[
g_t=\begin{pmatrix}e^t&0\\0&e^{-t}\end{pmatrix},
\]
and let \(r_{-\theta}\) rotate flat coordinates so that direction \(\theta\) becomes horizontal. The matrices act by postcomposition on flat coordinate charts. Under this forward-ray convention, define
\[
D_{\mathcal M}(X)=\left\{\theta\in\Theta:\ \text{for every compact }K\subseteq\mathcal M\text{ there exists }T\geq0\text{ such that }g_t r_{-\theta}X\notin K\text{ for all }t\geq T\right\}.
\]
The source record does not specify this convention or \(\mathcal M\); in particular, divergence in a fixed stratum and divergence after projection to the moduli space of underlying Riemann surfaces must not be silently identified.

\textbf{Question, under the stated conventions.} Can either of the following existence statements hold?
\[
\exists X\qquad 0<\dim_{\mathrm H}N(X)<\frac12,
\]
\[
\exists X\qquad 0<\dim_{\mathrm H}D_{\mathcal M}(X)<\frac12.
\]
Here Hausdorff dimension is computed in \(\Theta\), and each existential quantifier ranges over the intended class of translation surfaces, which the supplied record leaves unspecified. These are separate alternatives: their witnessing surfaces need not coincide. This interpretation of the source's word ``or'' requires confirmation; the record does not explicitly ask for the dimension of a union or require realization of every value in \((0,1/2)\).
\subsection{Short English statement}
Can some translation surface have a set of non-ergodic directions, or a set of divergent directions, whose Hausdorff dimension lies strictly between \(0\) and \(1/2\)? A precise final statement requires confirmation of the class of surfaces, the meaning of non-ergodicity, the ambient space and time convention for divergence, and the intended reading of ``or.''
\subsection{Sources}
\begin{itemize}
\item Local UnsolvedMath registry, permanent ID \(30001783\), supplied complete source record; exact source problem code \texttt{catalogue entry}. Collection: Oberwolfach Reports Open Problems. The record gives proposed year 2011 and names no individual proposer.
\item \emph{Billiards, Flat Surfaces, and Dynamics on Moduli Spaces} (2011), Oberwolfach Reports, pp.~1361--1427, DOI \texttt{10.4171/owr/2011/25}. Supplied source report and literature reference ``original OWR report'': \url{https://doi.org/10.4171/owr/2011/25}.
\item Supplied literature reference ``Cheung--Hubert--Masur dichotomy'': \url{https://arxiv.org/abs/0810.3956}.
\item Supplied literature reference ``2026 extension'': \url{https://doi.org/10.1017/etds.2025.10265}.
\item Literature assessment dated exactly \texttt{2026-08-21}: the supplied curation labels the problem open and states that the Cheung--Hubert--Masur dichotomy and a 2026 extension cover important slit-surface families, while realization of an intermediate Hausdorff dimension is not settled in general. This is supplied dated literature triage derived from the OWR curation, not an independently verified assessment; its time-sensitive status and cited claims require independent verification.
\item The registry labels the statement \texttt{corrected\_verified} and reports that its mathematical content was checked against the official Oberwolfach report, with notation, context, and question form normalized without changing the source-level claim. This is supplied verification metadata, not verification performed for this draft.
\end{itemize}
\end{document}
